\subsection{ROOT DIAGRAMS}

DEFINITION 2. A root diagram (or simply a diagram, if there is no risk of confusion) is a triple \(D = (M, M_{0}, R)\) where:

\((RD_{0})\) M is a free Z-module of finite type and the submodule \(M_{0}\) is a direct factor of M;

\((\mathrm{RD}_{\mathrm{I}})\) R is a finite subset of M; \(\mathrm{R} \cup \mathrm{M}_0\) generates the Q-vector space \(\mathbf{Q} \otimes \mathrm{M}\) ;

\((\mathrm{RD}_{\mathrm{II}})\) for all \(\alpha \in \mathrm{R}\) , there exists an element \(\alpha^{\vee}\) of \(\mathrm{M}^{*} = \mathrm{Hom}_{\mathbf{Z}}(\mathrm{M}, \mathbf{Z})\) such that \(\alpha^{\vee}(\mathrm{M}_{0}) = 0\) , \(\alpha^{\vee}(\alpha) = 2\) and the endomorphism \(x \mapsto x - \alpha^{\vee}(x)\alpha\) of M leaves R stable.

By Chap. VI, §1, no. 1, for all \(\alpha \in \mathbb{R}\) the element \(\alpha^{\vee}\) of \(\mathrm{M}^*\) is uniquely determined by \(\alpha\) ; we denote by \(s_{\alpha}\) the endomorphism \(x \mapsto x - \alpha^{\vee}(x)\alpha\) of M. Moreover (loc. cit.), the \(\mathbf{Q}\) -vector space \(\mathbf{Q} \otimes \mathrm{M}\) is the direct sum of \(\mathbf{Q} \otimes \mathrm{M}_0\) and the vector subspace V(R) generated by R, and R is a root system in V(R) (loc. cit., Def. 1).

The elements of R are called the roots of the root diagram D, and the elements \(\alpha^{\vee}\) of \(M^{*}\) the inverse roots. The group generated by the automorphisms \(s_{\alpha}\) of M is called the Weyl group of D and is denoted by \(\mathrm{W(D)}\) ; the elements of \(\mathrm{W(D)}\) induce the identity on \(M_{0}\) , and induce on \(\mathrm{V(R)}\) the transformations of the Weyl group of the root system R.

Examples. 1) For every free \(\mathbf{Z}\) -module of finite type \(M\) , the triple \((M, M, \varnothing)\) is a root diagram.

\begin{enumerate}
\def\labelenumi{\arabic{enumi})}
\setcounter{enumi}{1}
\tightlist
\item
  If \(D = (M, M_{0}, R)\) is a root diagram, let \(M_{0}^{*}\) be the orthogonal complement of \(V(R)\) in \(M^{*}\) , and let \(R^{\vee}\) be the set of inverse roots of D. Then \(D^{\vee} = (M^{*}, M_{0}^{*}, R^{\vee})\) is a root diagram, called the inverse of D. For all \(\alpha \in R\) , the symmetry \(s_{\alpha^{\vee}}\) of \(M^{*}\) is the contragredient automorphism of the symmetry \(s_{\alpha}\) of M; the map \(w \mapsto {}^{t} w^{-1}\) is an isomorphism from \(W(D)\) to \(W(D^{\vee})\) . Moreover, \(V(R^{\vee})\) can be naturally identified with the dual of the Q-vector space \(V(R)\) , \(R^{\vee}\) then being identified with the inverse root system of R.
\end{enumerate}

If the dual of \(\mathbf{M}^*\) is identified with M, the inverse diagram of \(\mathrm{D}^{\vee}\) is identified with D.

\begin{enumerate}
\def\labelenumi{\arabic{enumi})}
\setcounter{enumi}{2}
\item
  Let \((\mathfrak{g},\mathfrak{h})\) be a split reductive \(\mathbf{Q}\) -Lie algebra, and \(\mathrm{M} \subset \mathfrak{h}\) a permissible lattice (Chap. VIII, §12, no. 6, Def. 1). Let \(\mathrm{M}_0\) be the subgroup of M orthogonal to the roots of \((\mathfrak{g},\mathfrak{h})\) and \(\mathrm{R}^\vee\) the set of the \(H_{\alpha}\) , \(\alpha \in \mathrm{R}(\mathfrak{g},\mathfrak{h})\) . Then \((\mathrm{M},\mathrm{M}_0,\mathrm{R}^\vee)\) is a root diagram, and \((\mathrm{M}^*,\mathrm{M}_0^*,\mathrm{R}(\mathfrak{g},\mathfrak{h}))\) is the inverse diagram.
\item
  Let V be a vector space over Q and R a root system in V; denote by P(R) the group of weights of R and by Q(R) the group of radical weights of R (Chap. VI, §1, no. 9). Then (Q(R), 0, R) and (P(R), 0, R) are root diagrams. A root diagram (M, M₀, S) is isomorphic to a diagram of the form (Q(R), 0, R) (resp. (P(R), 0, R)) if and only if M is generated by S (resp. M* is generated by S \(^{v}\) ).
\end{enumerate}

For every subgroup X of P(R) containing Q(R), (X, 0, R) is a root diagram and, up to isomorphism, every diagram (M, M₀, S) such that M₀ = 0, in other words such that S generates a subgroup of M of finite index, arises in this way.

The root diagram \((\mathrm{M},\mathrm{M}_0,\mathrm{R})\) is said to be reduced if the root system \(\mathbf{R}\) is reduced (in other words (Chap. VI, §1, no. 4) if the relations \(\alpha ,\beta \in \mathbb{R}\) , \(\lambda \in \mathbf{Z}\) , \(\beta = \lambda \alpha\) imply that \(\lambda = 1\) or \(\lambda = -1\) ). The diagrams in Examples 1) and 3) are reduced.

